\section{Direct Portfolio Learning}
\label{sec:direct_learning}
To estimate the portfolio policy we use a kernel representation \citet{filipovicschneider2024,nagel2025} because it provides a common framework for
portfolio policies of very different forms, as shown in table .... 
Indeed, depending on the kernel, the same estimator can generate static, linear,
polynomial or sieve-based, and flexible nonlinear mappings, generalizing respectively the static mean--variance portfolio of \citet{brittenjones1999}, the linear parametric portfolio
policies of \citet{brandtsantaclaravalkanov2009}, and the semiparametric
spline-based policies of \citet{caldeirasantostorrent2023}.\\
Kernel methods are also closely connected to the recent work of
\citet{nagel2025,kellymalamud2025}. In particular, random Fourier features
\citep{rahimi2007random} provide finite-dimensional approximations to
stationary kernels, while the corresponding kernel estimator can be viewed as
the infinite-feature limit. In this sense, the kernel formulation provides a
natural benchmark for portfolio models whose nominal representation complexity
is allowed to become arbitrarily large.

Kernel methods are also closely related to neural-network representations.
A random-feature approximation can be interpreted as a neural network with a
fixed, randomly generated hidden layer, in which only the final linear
coefficients are estimated. More generally, neural tangent kernels arise as
the infinite-width limits of particular neural-network training regimes
\citep{jacot2018neural}.

Let
\[
K:\mathscr Z\times\mathscr Z\rightarrow\mathbb R
\]
be a positive-definite scalar kernel and let $\Hh_K$ denote the corresponding
space of characteristic-based portfolio policies.

\begin{assumption}[Portfolio representation]
\label{ass:kernel}
The space $\mathscr Z$ is Polish. Fix a jointly measurable kernel $K$,
with separable RKHS $\Hh_K$, satisfying
\[
\sup_{z\in\mathscr Z}K(z,z)\leq\kappa_K<\infty.
\]
The representation is fixed across economies and sample sizes.
\end{assumption}
Let $W_K^\star$ denote the minimum-norm population minimizer in $\Hh_K$,
when attained. Write $W^\star_\varepsilon=W^\star_{\varepsilon,K}$ and
$\Sigma_{\Hh_K,\varepsilon}$ when distinguishing economies.\\
The empirical regularized counterpart of
the population problem in Equation~\eqref{eq:population_loss} is
\begin{equation}
\begin{aligned}
\widehat W_{\lambda,T}
&:=
\argmin_{W\in\Hh_K}
\left\{
\frac1T\sum_{t=1}^T
\left(
1-\frac1N\sum_{i=1}^N
W(Z_{i,t})R_{i,t+1}
\right)^2
+
\lambda\|W\|_{\Hh_K}^2
\right\}
\\[0.5em]
&=
\argmin_{W\in\Hh_K}
\left\{
\widehat Q(W)
+
\lambda\|W\|_{\Hh_K}^2
\right\}
=
\min_{\substack{U\in\Hh\\ \|U\|_{\Hh}=1}}
\left\{
\frac{1}{
1+\widehat{\SR}^2_{\lambda,T}
\left(R^p_{t+1}(U)\right))\!
}
\right\},
\qquad \lambda>0.
\end{aligned}
\label{eq:estimator}
\end{equation}
where
\[
\widehat{\SR}_{\lambda,T}
\left(R^p_{t+1}(U)\right)
:=
\frac{
\frac1T
\sum_{t=1}^T
R^p_{t+1}(U)
}{
\sqrt{
\frac1T
\sum_{t=1}^T
\left(
R^p_{t+1}(U)
-
\frac1T
\sum_{s=1}^T
R^p_{s+1}(U)
\right)^2
+\lambda
}
}.
\]
The economic structure of the estimator becomes transparent by viewing
kernel-managed portfolio payoffs as elements of the portfolio policy space.
At each date $t$, define
\[
X_t
:=
\frac{1}{N}
\sum_{i=1}^{N}
R_{i,t+1}K(\cdot,Z_{i,t})
\in\Hh_K.
\]
Thus $X_t$ is a random element of $\Hh_K$. If evaluated at
$z\in\mathscr Z$, it gives
\[
X_t(z)
=
\frac{1}{N}
\sum_{i=1}^{N}
R_{i,t+1}K(z,Z_{i,t}) \in R.
\]
the payoff generated by the kernel-managed portfolio that assigns stock $i$
the weight $K(z,Z_{i,t})/N$.\\
Its population mean is
\[
\mu_{\Hh_K}
:=
\E[X_t]
\in\Hh_K,
\]
while its population second-moment operator is
\[
\Sigma_{\Hh_K}
:=
\E[X_t\otimes X_t]
:
\Hh_K\rightarrow\Hh_K,
\]
where, for every $V\in\Hh_K$,
\[
(X_t\otimes X_t)V
=
\langle V,X_t\rangle_{\Hh_K}X_t.
\]
The first moment captures expected managed-portfolio payoffs, while the
second-moment operator describes the magnitude of realized payoffs across
portfolio directions.\\
Their sample counterparts are the element
\[
\widehat\mu_T
=
\frac{1}{T}\sum_{t=1}^{T}X_t
\in\Hh_K,
\]
and the operator
\[
\widehat\Sigma_T
=
\frac{1}{T}\sum_{t=1}^{T}X_t\otimes X_t
:
\Hh_K\rightarrow\Hh_K.
\]

\Needspace{8\baselineskip}
\begin{proposition}[Managed-portfolio Markowitz representation]
\label{prop:managed_markowitz}

The population-optimal minimum-norm policy satisfies
\begin{equation}
\mu_{\Hh_K}
=
\Sigma_{\Hh_K} W_K^\star.
\label{eq:population_markowitz_representation}
\end{equation}
Equivalently,
\[
W_K^\star
=
\Sigma_{\Hh_K}^{\dagger}\mu_{\Hh_K},
\]
where $\Sigma_{\Hh_K}^{\dagger}$ denotes the pseudo inverse.\\
The regularized sample policy in Equation~\eqref{eq:estimator} satisfies
\begin{equation}
\widehat W_{\lambda,T}
=
\left(
\widehat\Sigma_T+\lambda I
\right)^{-1}
\widehat\mu_T,
\qquad
\lambda>0.
\label{eq:markowitz_representation}
\end{equation}

\end{proposition}

Equation~\eqref{eq:markowitz_representation} provides the structural link to
classical portfolio choice. The investor is still combining expected returns
against second moments, as in Markowitz. The difference is that the investment
objects are now characteristic-managed portfolios weighted with the kernels.

Regularization has a familiar portfolio interpretation. As emphasized by
\citet{kozaknagelsantosh2020}, portfolio directions with very little payoff
variation but non-negligible expected returns resemble near-arbitrage
opportunities and should therefore be treated with particular caution.
Regularization has a familiar portfolio interpretation. As emphasized by
\citet{kozaknagelsantosh2020}, weak portfolio directions are particularly
dangerous because small estimation errors can translate into very large
portfolio positions. In their eigenportfolio setting, a low-variance direction
with a sufficiently large estimated expected return can even resemble a
near-arbitrage opportunity.\\
The same instability appears in our managed-portfolio representation. Let
$\Sigma_{\Hh_K}e_j=\mu_j e_j$, and let
$m_j:=\langle\mu_{\Hh_K},e_j\rangle_{\Hh_K}$.
Without regularization, the portfolio weight along direction $e_j$ is
proportional to
\[
\widehat W_{0,j}
=
\frac{\widehat m_j}{\widehat\mu_j}.
\]
Hence, when $\widehat\mu_j$ is small, even modest estimation error in
$\widehat m_j$ can generate a very large portfolio position. Regularization
replaces the small eigenvalue in the denominator by
$\widehat\mu_j+\lambda$:
\[
\widehat W_{\lambda,j}
=
\frac{\widehat m_j}{\widehat\mu_j+\lambda}
=
\frac{\widehat\mu_j}{\widehat\mu_j+\lambda}
\widehat W_{0,j}.
\]
Thus, $\widehat\mu_j/(\widehat\mu_j+\lambda)$ is the spectral shrinkage factor.


The spectral filter derived above suggests a natural measure of how much of the
managed-portfolio opportunity set is effectively used by the investor. The
relevant notion of portfolio complexity is not the number of characteristics,
parameters, or random features in the representation, but the number of
economically distinct portfolio directions that remain active after shrinkage.

This motivates the following definition.

\begin{definition}[Effective portfolio complexity]
\label{def:effective_complexity}
For a fixed portfolio representation $K$ and $\lambda>0$, define
\begin{equation}
\mathcal C_K(\lambda)
:=
\operatorname{Tr}\!\left[
\Sigma_{\Hh_K}
(\Sigma_{\Hh_K}+\lambda I)^{-1}
\right]
=
\sum_{j\geq1}
\frac{\mu_j}{\mu_j+\lambda}.
\label{eq:effective_complexity}
\end{equation}
Effective portfolio complexity is the shrinkage-adjusted number of
managed-portfolio directions used by the portfolio rule.
\end{definition}

The definition follows directly from the spectral shrinkage mechanism already
introduced. Each managed-portfolio direction contributes between zero and one
to $\mathcal C_K(\lambda)$. If $\mu_j\gg\lambda$, its shrinkage factor is close
to one and the direction is effectively active. If $\mu_j\ll\lambda$, the
direction contributes little and is effectively suppressed.

Effective complexity therefore measures how deeply into the
managed-portfolio spectrum the investor is actually operating.

This notion is fundamentally different from nominal representation size. A
model may contain thousands, or infinitely many, basis directions while the
portfolio decision itself remains effectively low dimensional because
regularization suppresses most weak directions. Conversely, a rich
representation can be valuable because it makes additional investment
opportunities available once sufficient market history exists to identify them.

